\chapter{Conclusions and Future Scope for Improvement}

\section{Conclusions}

This project successfully demonstrates that Zero-Knowledge Proof based authentication is not merely a theoretical concept but a practical, deployable solution for real-world privacy-preserving identity verification. ZeroAuth was designed, implemented, and deployed as a fully functional three-component system --- a React Native mobile Wallet, a Node.js Relay server, and a JavaScript Web SDK --- each carefully engineered to work in concert while ensuring that sensitive credential data never leaves the user's device at any point in the verification process.

The Wallet application enables users to store multiple types of credentials locally and generate Groth16 zk-SNARK proofs using Circom circuits compiled to WebAssembly, executed within a sandboxed WebView bridge. The Relay server manages verification sessions with a five-minute TTL, validates incoming proof structures, performs server-side cryptographic verification using snarkjs, and enforces proof replay protection through SHA-256 hashing --- all backed by a Supabase PostgreSQL database with Row-Level Security policies. The Web SDK reduces integration effort for verifier websites to a single component or script tag, making the technology accessible to developers without deep cryptographic knowledge.

Key achievements include privacy by design with no private attributes ever transmitted, multi-layer cryptographic security combining Poseidon commitments, Ed25519 identity, Groth16 proofs, and SHA-256 replay protection, standards compliance with the W3C did:key specification, and full production deployment on cloud infrastructure with live demo sites. The system proved functionally correct across all tested scenarios: valid proofs are accepted, tampered or replayed proofs are rejected, expired sessions are cleaned up automatically, and the mobile wallet operates correctly on Android devices ranging from high-end to low-end hardware. ZeroAuth demonstrates that privacy-preserving authentication can be both cryptographically rigorous and practically usable, paving the way for broader adoption of Zero-Knowledge Proofs in identity systems.

\newpage
\section{Future Scope for Improvement}

While ZeroAuth achieves its intended objectives, several directions exist for future enhancement. On the cryptographic front, migrating from Groth16 to a universal zk-SNARK system such as PLONK or Halo2 would eliminate the per-circuit trusted setup ceremony, making the system more scalable and operationally simpler. Deploying auto-generated Solidity verifier contracts on Ethereum or Polygon would enable fully decentralised, trustless proof verification without a centralised Relay server. Integrating BBS+ signature-based Verifiable Credentials would allow issuer-signed credentials and true selective disclosure of multiple attributes from a single credential.

On the platform side, publishing the Wallet to the Apple App Store and implementing a full credential issuance SDK for institutions such as universities and employers would complete the identity ecosystem. Performance improvements --- such as replacing the WebView bridge with a native WASM module and adding WebSocket push notifications instead of polling --- would further enhance the user experience. Finally, binding wallet credentials to device biometrics, implementing comprehensive audit logging, and adding a decentralised credential revocation registry would strengthen the security posture of the system for production-scale deployments.

